% GENERATED FILE. Edit canonical lessons, then run npm run generate:books.
% canonical-source-hash: fnv1a64:9691b9528b135962
% canonical-lessons: TA-C332-por, TA-C332-patai, TA-C332-pataiyetu, TA-C332-takku, TA-C332-katarpatai, TA-R332-review-326-332

\chapter{War}
\label{ch:ta-war}
\hlchaptermodality{332}

\begin{chapteropening}
\textbf{By the end of this chapter:} \emph{I can talk about wars and armies.}

Complete the last lesson of chapter 332: I can talk about wars and armies.
\end{chapteropening}

\section[\texorpdfstring{\ta{போர்} (pōr)}{போர் (pōr)}]{\ta{போர்} (pōr) — a war; a battle}

\label{lesson:TA-C332-por}



\begin{quote}
Before the new one: say the Tamil for a donation, then the Tamil for to prosper.
\end{quote}

\subsection*{You'll want to know: \ta{போர்}}
\textbf{\ta{போர்}} — \emph{pōr} — \textquotedblleft{}a war; a battle\textquotedblright{}.

\begin{quote}
\ta{அந்தப்} \ta{போர்} \ta{பத்து} \ta{ஆண்டுகள்} \ta{நடந்தது}.
\end{quote}

\begin{quote}
\emph{antap pōr pattu āṇṭukaḷ naṭantatu.}
\end{quote}

\begin{quote}
That war went on for ten years.
\end{quote}

\ta{போர்} \ta{புரி}, to wage war.

\begin{cousinweb}
\textbf{\ta{போர்}} (\emph{pōr}) is native Dravidian: DEDR 4540 gives it under \ta{பொரு}, \textquotedblleft{}to fight, engage in battle\textquotedblright{}.
\end{cousinweb}

\begin{grammarlens}[title={What you've built}]
One more word for telling what happened.
\end{grammarlens}

\subsection*{Your turn}
\begin{itemize}
\raggedright
  \item \emph{Say it:} \emph{pōr}
  \item \emph{Say it:} \emph{pōr}, once more
  \item \emph{Say it:} \emph{antap pōr pattu āṇṭukaḷ naṭantatu.} — That war went on for ten years.
  \item \emph{Recall:} say the Tamil for a donation, then the Tamil for to prosper, then say \emph{pōr} again
\end{itemize}

\begin{tcolorbox}[breakable,colback=teal!4,colframe=teal!35!black,title={Before you move on}]
What does \ta{போர்} mean? (\textquotedblleft{}A war; a battle\textquotedblright{}.) Say it once more.
\end{tcolorbox}
\section[\texorpdfstring{\ta{படை} (paṭai)}{படை (paṭai)}]{\ta{படை} (paṭai) — an army; a force}

\label{lesson:TA-C332-patai}



\begin{quote}
Before the new one: say the Tamil for to prosper, then the Tamil for a war.
\end{quote}

\subsection*{You'll want to know: \ta{படை}}
\textbf{\ta{படை}} — \emph{paṭai} — \textquotedblleft{}an army; a force\textquotedblright{}.

\begin{quote}
\ta{மன்னரின்} \ta{படை} \ta{ஆற்றைக்} \ta{கடந்தது}.
\end{quote}

\begin{quote}
\emph{maṉṉariṉ paṭai āṟṟaik kaṭantatu.}
\end{quote}

\begin{quote}
The king's army crossed the river.
\end{quote}

Also a band of people.

\begin{cousinweb}
\textbf{\ta{படை}} (\emph{paṭai}) is native Dravidian (DEDR 3860): army, crowd, weapons, battle.
\end{cousinweb}

\begin{grammarlens}[title={What you've built}]
One more word for telling what happened.
\end{grammarlens}

\subsection*{Your turn}
\begin{itemize}
\raggedright
  \item \emph{Say it:} \emph{paṭai}
  \item \emph{Say it:} \emph{paṭai}, once more
  \item \emph{Say it:} \emph{maṉṉariṉ paṭai āṟṟaik kaṭantatu.} — The king's army crossed the river.
  \item \emph{Recall:} say the Tamil for to prosper, then the Tamil for a war, then say \emph{paṭai} again
\end{itemize}

\begin{tcolorbox}[breakable,colback=teal!4,colframe=teal!35!black,title={Before you move on}]
What does \ta{படை} mean? (\textquotedblleft{}An army; a force\textquotedblright{}.) Say it once more.
\end{tcolorbox}
\section[\texorpdfstring{\ta{படையெடு} (paṭaiyeṭu)}{படையெடு (paṭaiyeṭu)}]{\ta{படையெடு} (paṭaiyeṭu) — to march against; to invade}

\label{lesson:TA-C332-pataiyetu}



\begin{quote}
Before the new one: say the Tamil for a war, then the Tamil for an army.
\end{quote}

\subsection*{You'll want to know: \ta{படையெடு}}
\textbf{\ta{படையெடு}} — \emph{paṭaiyeṭu} — \textquotedblleft{}to march against; to invade\textquotedblright{}.

\begin{quote}
\ta{பக்கத்து} \ta{நாட்டு} \ta{மன்னர்} \ta{நம்} \ta{நாட்டின்} \ta{மீது} \ta{படையெடுத்தார்}.
\end{quote}

\begin{quote}
\emph{pakkattu nāṭṭu maṉṉar nam nāṭṭiṉ mītu paṭaiyeṭuttār.}
\end{quote}

\begin{quote}
The neighbouring king marched against our country.
\end{quote}

Past (I …): \ta{படையெடுத்தேன்}, \emph{paṭaiyeṭuttēṉ}.

The target takes \ta{மீது}: \ta{நாட்டின்} \ta{மீது} \ta{படையெடு}.

\begin{cousinweb}
\textbf{\ta{படையெடு}} (\emph{paṭaiyeṭu}) is a Tamil compound built from \ta{படை}, \textquotedblleft{}army\textquotedblright{} (DEDR 3860), and \ta{எடு}, \textquotedblleft{}to take up, raise\textquotedblright{} (DEDR 851): to raise an army against.
\end{cousinweb}

\begin{grammarlens}[title={What you've built}]
One more word for telling what happened.
\end{grammarlens}

\subsection*{Your turn}
\begin{itemize}
\raggedright
  \item \emph{Say it:} \emph{paṭaiyeṭu}
  \item \emph{Say it:} \emph{paṭaiyeṭu}, once more
  \item \emph{Say it:} \emph{pakkattu nāṭṭu maṉṉar nam nāṭṭiṉ mītu paṭaiyeṭuttār.} — The neighbouring king marched against our country.
  \item \emph{Recall:} say the Tamil for a war, then the Tamil for an army, then say \emph{paṭaiyeṭu} again
\end{itemize}

\begin{tcolorbox}[breakable,colback=teal!4,colframe=teal!35!black,title={Before you move on}]
What does \ta{படையெடு} mean? (\textquotedblleft{}To march against; to invade\textquotedblright{}.) Say it once more.
\end{tcolorbox}
\section[\texorpdfstring{\ta{தாக்கு} (tākku)}{தாக்கு (tākku)}]{\ta{தாக்கு} (tākku) — to attack; to strike}

\label{lesson:TA-C332-takku}



\begin{quote}
Before the new one: say the Tamil for an army, then the Tamil for to march against.
\end{quote}

\subsection*{You'll want to know: \ta{தாக்கு}}
\textbf{\ta{தாக்கு}} — \emph{tākku} — \textquotedblleft{}to attack; to strike\textquotedblright{}.

\begin{quote}
\ta{படைவீரர்கள்} \ta{இரவில்} \ta{கோட்டையைத்} \ta{தாக்கினார்கள்}.
\end{quote}

\begin{quote}
\emph{paṭaivīrarkaḷ iravil kōṭṭaiyait tākkiṉārkaḷ.}
\end{quote}

\begin{quote}
The soldiers attacked the fort at night.
\end{quote}

Past (I …): \ta{தாக்கினேன்}, \emph{tākkiṉēṉ}.

Also of illness or weather: \ta{காய்ச்சல்} \ta{அவனைத்} \ta{தாக்கியது}.

\begin{cousinweb}
\textbf{\ta{தாக்கு}} (\emph{tākku}) is native Dravidian (DEDR 3150): to collide, strike, attack.
\end{cousinweb}

\begin{grammarlens}[title={What you've built}]
One more word for telling what happened.
\end{grammarlens}

\subsection*{Your turn}
\begin{itemize}
\raggedright
  \item \emph{Say it:} \emph{tākku}
  \item \emph{Say it:} \emph{tākku}, once more
  \item \emph{Say it:} \emph{paṭaivīrarkaḷ iravil kōṭṭaiyait tākkiṉārkaḷ.} — The soldiers attacked the fort at night.
  \item \emph{Recall:} say the Tamil for an army, then the Tamil for to march against, then say \emph{tākku} again
\end{itemize}

\begin{tcolorbox}[breakable,colback=teal!4,colframe=teal!35!black,title={Before you move on}]
What does \ta{தாக்கு} mean? (\textquotedblleft{}To attack; to strike\textquotedblright{}.) Say it once more.
\end{tcolorbox}
\section[\texorpdfstring{\ta{கடற்படை} (kaṭaṟpaṭai)}{கடற்படை (kaṭaṟpaṭai)}]{\ta{கடற்படை} (kaṭaṟpaṭai) — a navy}

\label{lesson:TA-C332-katarpatai}



\begin{quote}
Before the new one: say the Tamil for to march against, then the Tamil for to attack.
\end{quote}

\subsection*{You'll want to know: \ta{கடற்படை}}
\textbf{\ta{கடற்படை}} — \emph{kaṭaṟpaṭai} — \textquotedblleft{}a navy\textquotedblright{}.

\begin{quote}
\ta{அந்த} \ta{மன்னரிடம்} \ta{பெரிய} \ta{கடற்படை} \ta{இருந்தது}.
\end{quote}

\begin{quote}
\emph{anta maṉṉariṭam periya kaṭaṟpaṭai iruntatu.}
\end{quote}

\begin{quote}
That king had a large navy.
\end{quote}

\ta{கடல்}, sea + \ta{படை}, army; \ta{ல்} turns to \ta{ற்} before \ta{ப}.

\begin{cousinweb}
\textbf{\ta{கடற்படை}} (\emph{kaṭaṟpaṭai}) is a Tamil compound built from \ta{கடல்}, \textquotedblleft{}sea\textquotedblright{} (DEDR 1118), and \ta{படை}, \textquotedblleft{}army\textquotedblright{} (DEDR 3860).
\end{cousinweb}

\begin{grammarlens}[title={What you've built}]
One more word for telling what happened.
\end{grammarlens}

\subsection*{Your turn}
\begin{itemize}
\raggedright
  \item \emph{Say it:} \emph{kaṭaṟpaṭai}
  \item \emph{Say it:} \emph{kaṭaṟpaṭai}, once more
  \item \emph{Say it:} \emph{anta maṉṉariṭam periya kaṭaṟpaṭai iruntatu.} — That king had a large navy.
  \item \emph{Recall:} say the Tamil for to march against, then the Tamil for to attack, then say \emph{kaṭaṟpaṭai} again
\end{itemize}

\begin{tcolorbox}[breakable,colback=teal!4,colframe=teal!35!black,title={Before you move on}]
What does \ta{கடற்படை} mean? (\textquotedblleft{}A navy\textquotedblright{}.) Say it once more.
\end{tcolorbox}
\section[(dialogue)]{Looking back: from The Past to War}

\label{lesson:TA-R332-review-326-332}



\begin{quote}
Say the Tamil for history and the past.
\end{quote}

\subsection*{Your turn}
Say these aloud:

\begin{itemize}
\raggedright
  \item history, the past, ancient times, a period, a decade
  \item ancestors, a dynasty, to settle, a long time, ancestral roots
  \item a king, a queen, a prince, a princess, a throne
  \item a royal court, a crown, to crown a king, a reign, an envoy
  \item an empire, an emperor, a treasury, a royal tribute, to flourish
  \item a royal decree, a copper-plate grant, a patron, a donation, to prosper
  \item a war, an army, to march against, to attack, a navy
\end{itemize}

\begin{tcolorbox}[breakable,colback=teal!4,colframe=teal!35!black,title={Before you move on}]
History? (\textbf{\ta{வரலாறு}}, \emph{varalāṟu}.) The past? (\textbf{\ta{கடந்த} \ta{காலம்}}, \emph{kaṭanta kālam}.) Ancient times? (\textbf{\ta{பழங்காலம்}}, \emph{paḻaṅkālam}.) A period? (\textbf{\ta{காலகட்டம்}}, \emph{kālakaṭṭam}.) A decade? (\textbf{\ta{தசாப்தம்}}, \emph{tacāptam}.) Ancestors? (\textbf{\ta{முன்னோர்}}, \emph{muṉṉōr}.) A dynasty? (\textbf{\ta{வம்சம்}}, \emph{vamcam}.) To settle? (\textbf{\ta{குடியேறு}}, \emph{kuṭiyēṟu}.) A long time? (\textbf{\ta{நெடுங்காலம்}}, \emph{neṭuṅkālam}.) Ancestral roots? (\textbf{\ta{பூர்வீகம்}}, \emph{pūrvīkam}.) A king? (\textbf{\ta{மன்னர்}}, \emph{maṉṉar}.) A queen? (\textbf{\ta{ராணி}}, \emph{rāṇi}.) A prince? (\textbf{\ta{இளவரசன்}}, \emph{iḷavaracaṉ}.) A princess? (\textbf{\ta{இளவரசி}}, \emph{iḷavaraci}.) A throne? (\textbf{\ta{அரியணை}}, \emph{ariyaṇai}.) A royal court? (\textbf{\ta{அரசவை}}, \emph{aracavai}.) A crown? (\textbf{\ta{மகுடம்}}, \emph{makuṭam}.) To crown a king? (\textbf{\ta{முடிசூட்டு}}, \emph{muṭicūṭṭu}.) A reign? (\textbf{\ta{ஆட்சி}}, \emph{āṭci}.) An envoy? (\textbf{\ta{தூதர்}}, \emph{tūtar}.) An empire? (\textbf{\ta{பேரரசு}}, \emph{pēraracu}.) An emperor? (\textbf{\ta{சக்கரவர்த்தி}}, \emph{cakkaravartti}.) A treasury? (\textbf{\ta{கருவூலம்}}, \emph{karuvūlam}.) A royal tribute? (\textbf{\ta{கப்பம்}}, \emph{kappam}.) To flourish? (\textbf{\ta{ஓங்கு}}, \emph{ōṅku}.) A royal decree? (\textbf{\ta{அரசாணை}}, \emph{aracāṇai}.) A copper-plate grant? (\textbf{\ta{செப்பேடு}}, \emph{ceppēṭu}.) A patron? (\textbf{\ta{வள்ளல்}}, \emph{vaḷḷal}.) A donation? (\textbf{\ta{கொடை}}, \emph{koṭai}.) To prosper? (\textbf{\ta{செழி}}, \emph{ceḻi}.) A war? (\textbf{\ta{போர்}}, \emph{pōr}.) An army? (\textbf{\ta{படை}}, \emph{paṭai}.) To march against? (\textbf{\ta{படையெடு}}, \emph{paṭaiyeṭu}.) To attack? (\textbf{\ta{தாக்கு}}, \emph{tākku}.) A navy? (\textbf{\ta{கடற்படை}}, \emph{kaṭaṟpaṭai}.)
\end{tcolorbox}
